% docs/texte/10_engineering/licence.tex -- chapter "Licence" (Interferosim)
% Version: 0.4.0 (2026-11-23): two licences -- code MIT, documentation,
%   own figures and notebook texts CC BY 4.0 (LICENSE-docs); overview
%   table, label ch:licence
%   0.3.0 (2026-11-03): chapter written out (scope, permissions,
%   third-party material, citation) for release 1.0.0
%   0.2.0: MIT instead of CC BY 4.0
% ================================================================
\index{licence}\index{MIT License}
\index{CC BY 4.0}\label{ch:licence}  % 0.4.0

% 0.4.0: overview of both licences
\section{Overview}

Interferosim uses two licences: the \textbf{MIT License} for the
software and the \textbf{Creative Commons Attribution~4.0
International} licence (CC~BY~4.0) for the documentation and the
figures. Both texts are in the project root.\vspace{-\baselineskip}

\begin{center}
\begin{tabularx}{\linewidth}{@{}LLl@{}}
\toprule
Part & Licence & File \\
\midrule
source code (\texttt{src/}), tests, configuration files, scripts, live
  demo (\texttt{demo/}), code cells of the notebooks
  & MIT License & \texttt{LICENSE} \\
documentation (\texttt{docs/}: \LaTeX{} and Sphinx), own figures
  (\texttt{docs/images/}, figures generated by the notebooks), text
  cells of the notebooks
  & CC~BY~4.0 & \texttt{LICENSE-docs} \\
third-party figures and data & licence of the source
  & caption, file header \\
\bottomrule
\end{tabularx}
\end{center}

\section{MIT License}

Since version~1.0.0 the software of Interferosim is released under the
\textbf{MIT License}; the full text is in \texttt{LICENSE} in the project root. The
licence is also stated in \texttt{pyproject.toml} and
\texttt{CITATION.cff}.

\begin{center}
\begin{tabular}{@{}ll@{}}
\toprule
Licence & MIT License \\
Copyright & \copyright~2026 Erich Schubert \\
Repository & \url{https://github.com/eschubi/interferosim} (public) \\
% 0.4.0: documentation no longer under MIT (now CC BY 4.0, see below)
%Applies to & source code, tests, notebooks, configuration files, \\
%           & live demo and documentation \\
Applies to & source code, tests, configuration files, scripts, \\
           & live demo, code cells of the notebooks \\
\bottomrule
\end{tabular}
\end{center}

The MIT License allows anyone to use, copy, modify, merge, publish,
distribute, sublicense and sell the software, also in teaching material
and in other projects. The only condition is that the copyright notice
and the licence text are kept in all copies or substantial portions of
the software. The software is provided \enquote{as is}, without warranty
of any kind.

Up to version~0.9.7 the repository was private and marked with the
Creative Commons Attribution~4.0 licence (CC~BY~4.0). For the public
release the licence was changed to the MIT License, which is common for
software and fits the permissive licences of the dependencies (numpy,
scipy, astropy and matplotlib use BSD-style licences).

% 0.4.0: CC BY 4.0 added again for documentation and figures
\section{CC BY 4.0 for documentation and figures}

The MIT License is written for software. For texts and images the
Creative Commons licences are the common choice (Creative Commons itself
advises against using them for software). Since 2026-11-23 the
documentation, the own figures and the text cells of the notebooks are
therefore released under CC~BY~4.0; the full text is in
\texttt{LICENSE-docs}.

CC~BY~4.0 allows anyone to share and adapt the material for any
purpose, also commercially, for example in lecture notes, slides and
theses. The condition is attribution: name the author (Erich Schubert,
Interferosim), give the licence with a link to
\url{https://creativecommons.org/licenses/by/4.0/} and indicate whether
changes were made. A figure generated by the notebooks is part of the
documentation in this sense; the code that generates it stays under the
MIT License.

\section{Third-party material}

% 0.4.0: "not covered by the MIT License" -> neither licence
Figures and data taken from other sources are not covered by either
licence of the project; their source is named in the caption or in the
header of the file. Figures redrawn after a published figure name the
original. Third-party images whose licence is not clear are not put
under version control (chapter \enquote{Version Control}).

\section{Citation}

If Interferosim is used in teaching material, theses or publications,
please cite it with the metadata in \texttt{CITATION.cff}; GitHub shows
them under \enquote{Cite this repository}.
